\section*{अध्याय \ref{ch:b2:magmatic-differentiation} --- मैग्मीय विभेदन}

\begin{solution}{exo:b2:magmatic-differentiation:1}
(क) An$_{50}$ का लिक्विडस \qty{1450}{\celsius} है। पहला क्रिस्टल उस तापमान पर
सॉलिडस पर, An$_{66.7}$ (1372) और An$_{83.3}$ (1465) के बीच, स्थित है:
$66.7 + 16.6\times(1450 - 1372)/(1465 - 1372) = 66.7 + 13.9 = \qty{80.6}{\percent}$
एनॉर्थाइट। (ख) \qty{1394}{\celsius} पर द्रव लिक्विडस पर है,
An$_{33.3}$; क्रिस्टल सॉलिडस पर, $66.7 + 16.6\times22/93 =
\qty{70.6}{\percent}$। उत्तोलक नियम से क्रिस्टलित अंश $(50 -
33.3)/(70.6 - 33.3) = 16.7/37.3 = 0.448$ है। (ग) अंतिम द्रव तब लुप्त होता है जब
क्रिस्टल An$_{50}$ तक पहुँचते हैं, सॉलिडस तापमान \qty{1287}{\celsius} पर; वहाँ
लिक्विडस An$_{11.1}$ (1265) और An$_{20}$ (1334) के बीच है: $11.1 +
8.9\times22/69 = \qty{13.9}{\percent}$ एनॉर्थाइट।
\end{solution}

\begin{solution}{exo:b2:magmatic-differentiation:2}
$\bar D_{\ce{Sr}} = 0.15\times0.000063 + 0.35\times0.157 + 0.50\times2.7 = 0.00001 +
0.0550 + 1.35 = 1.40$; $\bar D_{\ce{Zr}} = 0.0001 + 0.0683 + 0.0005 = 0.069$;
$\bar D_{\ce{Rb}} = 0.00001 + 0.0012 + 0.0125 = 0.014$। केवल स्ट्रॉन्शियम, जिसे
प्लैजियोक्लेज़ धारण करता है, संगत है ($\bar D > 1$); ज़िरकोनियम और रुबिडियम असंगत हैं,
रुबिडियम अधिक।
\end{solution}

\begin{solution}{exo:b2:magmatic-differentiation:3}
(क) $C_L/C_0 = 0.5^{D-1}$: Sr $0.5^{0.40} = 0.758$ (यह घटता है); Zr $0.5^{-0.931} =
1.907$; Rb $0.5^{-0.986} = 1.981$, सीमा 2 के निकट। (ख) Zr/Rb
$0.5^{0.069 - 0.014} = 0.5^{0.055} = 0.963$ गुना बदलता है, यानी \qty{4}{\percent} से कम, जबकि हर
तत्त्व लगभग दोगुना हो गया: जैसा उपप्रमेय कहता है, दो \omterm{def:b2:elements-in-the-earth:partition}{असंगत तत्त्वों} का अनुपात
जनक \omterm{def:b2:magma-emplacement:magma}{मैग्मा} का मान बनाए रखता है, और इसलिए स्रोत के बारे में बताता है,
क्रिस्टलन के बारे में नहीं।
\end{solution}

\begin{solution}{exo:b2:magmatic-differentiation:4}
(क) \ce{SiO2} $+5.03$, \ce{MgO} $-1.95$, \ce{K2O} $+0.82$ (प्रतिशत), और पोटाश
$1.67/0.85 = 1.96$ गुना हुआ। (ख) पूर्णतः \omterm{def:b2:elements-in-the-earth:partition}{असंगत तत्त्व} के लिए
$C_L/C_0 = 1/F$, अतः $F = 0.85/1.67 = 0.509$: \omterm{def:b2:magma-emplacement:magma}{मैग्मा} के आधे भाग को
क्रिस्टलित होना पड़ता। (ग) \qty{12.5}{g} प्रति \qty{100}{g} हटाने से $F \approx 0.88$ बचता है और
पोटाश केवल लगभग $1/0.88 = 1.14$ गुना बढ़ता: अकेला क्रिस्टलन इस वृद्धि की
व्याख्या नहीं कर सकता। पोटाश बाहर से आया, पोटाश से समृद्ध (\qty{3.6}{\percent})
ग्रेनाइटी ज़ेनोलिथों से, जिन्हें \omterm{def:b2:magma-emplacement:magma}{मैग्मा} ने आत्मसात किया।
\end{solution}

\begin{solution}{exo:b2:magmatic-differentiation:5}
(क) $F = 0.8$, $D = 5$। भिन्नात्मक: $0.8^{4} = 0.410$; साम्य:
$1/(5 + 0.8\times(1 - 5)) = 1/1.8 = 0.556$। (ख) $F = 0.1$, $D = 0.01$। भिन्नात्मक:
$0.1^{-0.99} = 9.77$; साम्य: $1/(0.01 + 0.1\times0.99) = 1/0.109 = 9.17$।
(ग) $F \to 0$ होने पर रैले नियम $0$ की ओर जाता है, $D > 1$ के लिए, और अनंत की ओर,
$D < 1$ के लिए, जबकि साम्य नियम $1/D$ की ओर जाता है (समग्र संघटन वाले क्रिस्टलों के साथ
साम्य में अंतिम द्रव)। केवल \omterm{def:b2:magmatic-differentiation:crystallisation}{भिन्नात्मक क्रिस्टलन} ही किसी
द्रव से \omterm{def:b2:elements-in-the-earth:partition}{संगत तत्त्व} को पूरी तरह निकाल सकता है।
\end{solution}

\begin{solution}{exo:b2:magmatic-differentiation:6}
(क) एक सूत्र-इकाई में 1.6 \ce{MgO}, 0.4 \ce{FeO} और एक \ce{SiO2} होते हैं: $1.6\times
40.30 = 64.48$, $0.4\times71.84 = 28.74$, 60.08, कुल \qty{153.30}{g}। द्रव्यमान-प्रतिशत:
\ce{MgO} \qty{42.06}{\percent}, \ce{FeO} \qty{18.75}{\percent},
\ce{SiO2} \qty{39.19}{\percent}। (ख) \ce{MgO} पर उत्तोलक नियम: $x = (10.0 - 8.0)/(42.06
- 8.0) = 2.0/34.06 = 0.0587$, यानी लगभग \qty{6}{\percent} ओलिवीन। (ग) द्रव की
सिलिका: $(49.0 - 0.0587\times39.19)/(1 - 0.0587) = 46.70/0.9413 =
\qty{49.61}{\percent}$। बेसाल्ट से कम सिलिका वाले ओलिवीन को हटाने से द्रव की
सिलिका बढ़ती है, पर केवल धीरे-धीरे।
\end{solution}

\begin{solution}{exo:b2:magmatic-differentiation:7}
(क) $F(0) = 0.768$; $F(1000) = 0.1091 - 0.5906 + 0.7678 = 0.286$; $F(2000) = 0.4364
- 1.1813 + 0.7678 = 0.023$। (ख) विविक्तकर $(5.9064\times10^{-4})^2 -
4\times1.091\times10^{-7}\times0.7678 = 3.4885\times10^{-7} - 3.3507\times10^{-7} =
1.379\times10^{-8}$ है, जिसका वर्गमूल $1.174\times10^{-4}$ है; छोटा मूल
$H = (5.9064 - 1.174)\times10^{-4}/(2.182\times10^{-7}) = \qty{2169}{m}$ है, यानी
सैंडविच क्षितिज। (ग) $1 - 0.768 = 0.232$: अनावृत निचले क्षेत्र का आधार बनने से पहले
\omterm{def:b2:magma-emplacement:magma}{मैग्मा} का लगभग एक चौथाई भाग क्रिस्टलित हो चुका था, उसके नीचे के छिपे क्षेत्र में
और दीवारों तथा छत के अनुदिश सीमांत श्रेणी में।
\end{solution}

\begin{solution}{exo:b2:magmatic-differentiation:8}
\omterm{def:b2:magmatic-differentiation:cumulate}{क्यूम्यूलेट} में $D\,C_L = D\,C_0F^{D-1}$ होता है, अतः दो स्तरों के बीच
$\ln(\mathrm{Ni}_2/\mathrm{Ni}_1) = (D - 1)\ln(F_2/F_1)$। (क) $D - 1 = \ln(54/177)/
\ln(0.29/0.70) = (-1.187)/(-0.881) = 1.35$, अतः $D \approx 2.35$ (
$F \ge 0.23$ वाले 66 \omterm{def:b2:magmatic-differentiation:cumulate}{क्यूम्यूलेटों} पर समाश्रयण 2.38 देता है)। (ख) $D - 1 = \ln(6.2/54)/\ln(0.17/0.29) = (-2.164)/(-0.534)
= 4.05$, अतः $D \approx 5.05$। (ग) समग्र $D$ तब बदलता है जब क्रिस्टलित होता
समुच्चय बदलता है (कोई नया खनिज जो निकेल ग्रहण करता है, या कोई जो हट जाता है), और
खनिज--गलित गुणांक स्वयं द्रव के ठंडा होने और संघटन बदलने के साथ बढ़ते हैं,
जैसा ओलिवीन के लिए मापा गया है; \omterm{def:b2:magmatic-differentiation:cumulate}{क्यूम्यूलेटों} में फँसा द्रव
अनुमान को धुँधला करता है।
\end{solution}

\begin{solution}{exo:b2:magmatic-differentiation:9}
(क) $100 + 25.4 - 2.9 - 9.6 = \qty{112.9}{g}$। (ख) \ce{SiO2} के ग्राम: $55.53 +
25.4\times0.7248 - 2.9\times0.3919 - 9.6\times0.5054 = 55.53 + 18.41 - 1.14 - 4.85 =
\qty{67.95}{g}$, अतः $67.95/112.9 = \qty{60.19}{\percent}$, 60.16 के मुक़ाबले। \ce{MgO}:
$4.66 + 0.20 - 1.22 = \qty{3.64}{g}$, \qty{3.22}{\percent}, 3.19 के मुक़ाबले। \ce{K2O}:
$0.88 + 0.92 = \qty{1.80}{g}$, \qty{1.60}{\percent}, 1.75 के मुक़ाबले। (ग) सिलिका और मैग्नीशिया, तीन में से दो
नियंत्रक, 0.03 के भीतर मेल खाते हैं; पोटाश, जो नियंत्रक के रूप में प्रयुक्त नहीं हुआ, 0.15 कम है, अपने मान का लगभग
दसवाँ भाग: आत्मसात पदार्थ में औसत ज़ेनोलिथ से अधिक पोटाश था,
या उसका पोटाश-समृद्ध गलित वरीयता से ग्रहण किया गया।
\end{solution}

\begin{solution}{exo:b2:magmatic-differentiation:10}
(क) $z = (0.2 + 1.5 - 1)/(0.2 - 1) = 0.7/(-0.8) = -0.875$; $F^{-z} = 0.7^{0.875} =
0.7319$; $C^*/C_0 = 0.2\times0.5/(0.2 - 1 + 1.5) = 0.1/0.7 = 0.1429$, अतः $C/C_0 =
0.7319 + 0.1429\times(1 - 0.7319) = 0.7319 + 0.0383 = 0.770$। (ख) रैले:
$0.7^{0.5} = 0.837$। AFC \omterm{def:b2:magma-emplacement:magma}{मैग्मा} में स्ट्रॉन्शियम दो कारणों से कम है: जो क्रिस्टलित करता है उसका $r$ गुना पाते हुए अपने द्रव्यमान का
\qty{30}{\percent} खोने के लिए
उसे 0.30 के स्थान पर $(1 - 0.7)/(1 - 0.2) = 0.375$ क्रिस्टलित करना पड़ता है, और जो स्ट्रॉन्शियम वह पाता है
वह आधी समृद्ध चट्टान से आता है। (ग) $(0.7319\times0.7035 + 0.0383\times0.7200)/0.7702
= (0.5149 + 0.0276)/0.7702 = 0.70432$: अनुपात 0.0008 बढ़ गया है, जबकि
अकेला क्रिस्टलन उसे 0.7035 पर ही छोड़ देता।
\end{solution}

\begin{solution}{exo:b2:magmatic-differentiation:11}
(क) $1.1\times(1100 - 200) + 250 = 990 + 250 = \qty{1240}{J}$ प्रति ग्राम। (ख)
$25.4\times1240 = \qty{31496}{J}$, विलकॉक्स के \qty{31350}{J} के निकट। (ग)
$31496/370 = \qty{85.1}{g}$ प्लैजियोक्लेज़। (घ) द्रव्यमान-संतुलन ने केवल $2.9 +
9.6 = \qty{12.5}{g}$ क्रिस्टल हटाए, जिनकी गुप्त ऊष्मा (लगभग \qty{4600}{J}) आवश्यकता के
छठे भाग से भी कम पूरी करती है; $r = 25.4/12.5 = 2.03$, यानी बने क्रिस्टलों से दोगुनी चट्टान आत्मसात हुई,
जबकि ऊष्मा-बजट $r$ को 1 से काफ़ी नीचे रखता है। फूटा हुआ
\omterm{def:b2:magma-emplacement:magma}{मैग्मा} स्वयं यह ऊष्मा नहीं दे सकता था: \omterm{def:b2:magma-emplacement:magma}{मैग्मा} का कोई बड़ा द्रव्यमान गहराई में,
दृष्टि से ओझल, क्रिस्टलित होकर ठंडा हुआ होगा।
\end{solution}

\begin{solution}{exo:b2:magmatic-differentiation:12}
(क) चरण 1, $F$ 1 से 0.8 तक, $D = 5$ के साथ: गुणक $0.8^{4} = 0.4096$। चरण 2,
0.8 से 0.5 तक, $D = 1.5$ के साथ: नियम शेष द्रव पर लागू होता है, अतः
गुणक $(0.5/0.8)^{0.5} = 0.7906$ है। निकेल: $250\times0.4096\times0.7906 =
\qty{81.0}{ppm}$। (ख) $250\times0.5^{4} = \qty{15.6}{ppm}$, पाँच गुना बहुत कम। (ग)
रैले नियम स्थिर $D$ मानता है; जब क्रिस्टलित होता समुच्चय बदलता है,
$D$ उछलता है, और नियम को परिवर्तन पर पहुँचे द्रव-संघटन से फिर से शुरू करना पड़ता है,
जहाँ $F$ उसी द्रव के सापेक्ष मापा जाता है।
\end{solution}

\begin{solution}{pb:b2:magmatic-differentiation:1}
\noindent\textbf{भाग I।}

\textbf{1.} $\bar D_{\ce{Ni}} = 0.20\times5.0 + 0.50\times0 + 0.30\times2.0 = 1.60$।

\textbf{2.} $\bar D_{\ce{Zr}} = 0.20\times0.0007 + 0.50\times0.001 + 0.30\times0.195 =
0.00014 + 0.0005 + 0.0585 = 0.059$।

\textbf{3.} ऊपरी क्षेत्र: $\bar D_{\ce{Ni}} = 0.45\times2.0 = 0.90$; $\bar D_{\ce{Zr}} =
0.55\times0.001 + 0.45\times0.195 = 0.088$।

\textbf{4.} निकेल निचले क्षेत्र में संगत (1.60) है और ओलिवीन के लुप्त होने पर थोड़ा
असंगत (0.90) हो जाता है; ज़िरकोनियम दोनों में असंगत है।

\textbf{5.} क्लाइनोपाइरॉक्सीन: यह $0.0585$ देता है, निचले क्षेत्र के कुल $0.059$ में से, यानी लगभग
सारा, क्योंकि यही एकमात्र खनिज है जिसका गुणांक शून्य से दूर है।

\textbf{6.} $(C_{\ce{Ni}}/C_{\ce{Zr}})$ गुणक $F^{D_{\ce{Ni}} - D_{\ce{Zr}}}$ से बदलता है;
हर अर्धन पर, $0.5^{1.60 - 0.059} = 0.5^{1.541} = 0.344$।

\noindent\textbf{भाग II।}

\textbf{7.} $1.60\times150 = \qty{240}{ppm}$।

\textbf{8.} द्रव में $190/1.60 = \qty{118.75}{ppm}$ था; $F = (118.75/150)^{1/0.6} =
0.7917^{1.667} = 0.68$।

\textbf{9.} $80\times0.68^{-0.941} = 80\times1.438 = \qty{115}{ppm}$।

\textbf{10.} भिन्नात्मक: $150\times0.68^{0.6} = 150\times0.793 = \qty{119.0}{ppm}$;
साम्य: $150/(1.6 + 0.68\times(1 - 1.6)) = 150\times0.839 = \qty{125.8}{ppm}$।
\omterm{def:b2:magma-emplacement:magma}{मैग्मा} का एक तिहाई भाग क्रिस्टलित होने के बाद अंतर केवल \qty{6}{\percent} है:
मध्यम रूप से \omterm{def:b2:elements-in-the-earth:partition}{संगत तत्त्व} के लिए दोनों नियम आरंभ में निकट रहते हैं।

\textbf{11.} $150\times0.55^{0.6} = \qty{104.8}{ppm}$।

\textbf{12.} ठीक नीचे: $1.60\times104.8 = \qty{167.7}{ppm}$; ठीक ऊपर:
$0.90\times104.8 = \qty{94.3}{ppm}$। \omterm{def:b2:magmatic-differentiation:cumulate}{क्यूम्यूलेटों} की निकेल-प्रोफ़ाइल उस स्तर पर
अचानक \qty{44}{\percent} गिरती है जहाँ ओलिवीन परतों से लुप्त होता है:
एक सीढ़ी, कोई चिकना वक्र नहीं।

\textbf{13.} द्रव: $104.8\times(0.40/0.55)^{0.90 - 1} = 104.8\times1.032 =
\qty{108.2}{ppm}$; \omterm{def:b2:magmatic-differentiation:cumulate}{क्यूम्यूलेट} $0.90\times108.2 = \qty{97.4}{ppm}$। $\bar D < 1$ के साथ
निकेल अब द्रव में और \omterm{def:b2:magmatic-differentiation:cumulate}{क्यूम्यूलेटों} में धीरे-धीरे बढ़ता है: प्रोफ़ाइल अपनी
सीढ़ी के बाद धीरे-धीरे चढ़ती है।

\noindent\textbf{भाग III।}

\textbf{14.} $x_{\mathrm{ol}}(m_{\mathrm{ol},j} - d_j) + x_{\mathrm{pl}}
(m_{\mathrm{pl},j} - d_j) + x_{\mathrm{cpx}}(m_{\mathrm{cpx},j} - d_j) = 100\,(p_j -
d_j)$।

\textbf{15.} \ce{MgO}: $37.16\,x_{\mathrm{ol}} - 4.9\,x_{\mathrm{pl}} +
11.1\,x_{\mathrm{cpx}} = 410$; \ce{CaO}: $-7.3\,x_{\mathrm{ol}} +
7.06\,x_{\mathrm{pl}} + 13.2\,x_{\mathrm{cpx}} = 420$; \ce{Al2O3}:
$-14.1\,x_{\mathrm{ol}} + 17.6\,x_{\mathrm{pl}} - 11.1\,x_{\mathrm{cpx}} = 190$।

\textbf{16.} (विलोपन या क्रेमर नियम से) हल करने पर: $x_{\mathrm{ol}} = 9.03$,
$x_{\mathrm{pl}} = 30.84$, $x_{\mathrm{cpx}} = 20.32$ ग्राम प्रति \qty{100}{g}।

\textbf{17.} क्रिस्टल $9.03 + 30.84 + 20.32 = \qty{60.19}{g}$; द्रव
\qty{39.81}{g}: $F = 0.398$।

\textbf{18.} $(4950 - 9.03\times39.19 - 30.84\times50.54 - 20.32\times51.5)/39.81 =
(4950 - 353.9 - 1558.7 - 1046.5)/39.81 = 1990.9/39.81 = \qty{50.0}{\percent}$, जैसा
\omterm{def:b2:magma-emplacement:bodies}{डाइक} में है: जो ऑक्साइड नियंत्रक के रूप में प्रयुक्त नहीं हुआ, वह पुनः प्राप्त हो जाता है।

\textbf{19.} हटाए गए क्रिस्टल $9.03/60.19 = \qty{15}{\percent}$ ओलिवीन,
\qty{51}{\percent} प्लैजियोक्लेज़ और \qty{34}{\percent} क्लाइनोपाइरॉक्सीन हैं, दोनों
क्षेत्रों के खनिज-संघटनों के बीच। वे $F = 1$ से 0.55 तक क्रिस्टलित होते निचले समुच्चय
($0.45\times20 = 9.0$, $0.45\times50 = 22.5$, $0.45\times30 = 13.5$) और उसके बाद
0.55 से 0.40 तक ऊपरी समुच्चय ($0.15\times55 = 8.25$, $0.15\times45 = 6.75$) से मेल खाते हैं: 9.0,
30.75 और 20.25 ग्राम।

\textbf{20.} $0.30/0.75 = 0.40$।

\noindent\textbf{भाग IV।}

\textbf{21.} $0.55^{0.059 - 1} = 0.55^{-0.941} = 1.755$; $80\times1.755 =
\qty{140.4}{ppm}$।

\textbf{22.} $188/140.4 = 1.339 = (F/0.55)^{0.088 - 1}$, अतः $F = 0.55\times1.339^{-1/
0.912} = 0.55\times0.726 = 0.399$।

\textbf{23.} $(188/80)^{1/(0.07 - 1)} = 2.35^{-1.075} = 0.399$। \omterm{def:b2:elements-in-the-earth:partition}{असंगत
तत्त्वों} के लिए $C_L/C_0 \approx 1/F$, छोटा $D$ जो भी हो: उत्तर
$D$ के विवरणों पर मुश्किल से निर्भर करता है, जिससे \omterm{def:b2:elements-in-the-earth:partition}{असंगत तत्त्व} $F$ के सर्वोत्तम मापक बनते हैं।

\textbf{24.} प्लैजियोक्लेज़, जिसमें लोहा नहीं होता, क्रिस्टलों का आधा भाग बनाता है, और
पाइरॉक्सीन तथा ओलिवीन मिलकर जितनी मैग्नीशिया धारण करते हैं उससे कम लोहा धारण करते हैं: द्रव
सिलिका पाए बिना लोहे में अधिक समृद्ध और मैग्नीशिया
में निर्धन होता जाता है। लगभग स्थिर सिलिका पर
यह लौह-संवर्धन अगले अध्याय की थोलिआइटी श्रेणी का चिह्न है।

\textbf{25.} शुद्ध क्यूम्यूलस क्रिस्टल $0.088\times188 = \qty{16.5}{ppm}$ धारण करते। फँसे द्रव के
द्रव्यमान-अंश $t$ के साथ, $16.5(1 - t) + 188\,t = 47$, अतः $t = (47 -
16.5)/(188 - 16.5) = 30.5/171.5 = 0.18$: लगभग \qty{18}{\percent} फँसा द्रव।

\textbf{26.} तीनों अनुमान, मुख्य तत्त्व (0.398), पोटाश (0.40) और
ज़िरकोनियम (0.399), मेल खाते हैं: जब सबसे ऊपरी परत क्रिस्टलित हुई, तब
मूल \omterm{def:b2:magma-emplacement:magma}{मैग्मा} का \textbf{\qty[text-series-to-math=true]{40}{\percent}} भाग
अब भी द्रव था।
\end{solution}
